\documentclass[11pt]{article}
\usepackage{icmlko}
\begin{document}

\title{마지막 손수 고른 상수가 진짜 봉우리에 있었다 --- 그리고 씨앗 스물이 세 번째로 속였다}
\author{월드모델 실험실}
\date{노트 518 · 2026년 8월 4일}
\maketitle

\begin{abstract}
챔피언에 남은 \emph{손수 고른} 상수는 \texttt{TRAINCAP}\{만화$:0.5$\}
하나였다. 노트 517 이 그것을 빼면 안 된다는 것까지는 보였지만 $0.5$ 라는
값 자체는 노트 508 에서 ``절반''이라고 정한 그대로였다. 양옆을 쟀다.

\begin{center}
\begin{tabular}{lrrrrr}
\hline
씨앗 $20$ · 현 챔피언 대비 & 판 & 날짜 통제 & 무리 안 & KR 만화 & 비게임 앱 \\
\hline
만화 $0.25$ & $-0.0000$ & $-0.0002$ & $-0.0000$ & $-0.0021$ & $+0.0044$ \\
만화 $0.75$ & $\mathbf{-0.0026}$ & $\mathbf{-0.0024}$ & $\mathbf{-0.0027}$ & $-0.0002$ & $-0.0026$ \\
$+$세계애니 $0.5$ & $\mathbf{-0.0013}$ & $\mathbf{-0.0014}$ & $-0.0009$ & $-0.0003$ & $+0.0022$ \\
\hline
\end{tabular}
\end{center}

$0.75$ 는 셋이 $0$ 밖 음수, 세계애니 추가는 둘이 $0$ 밖 음수 --- 둘 다
기각. \textbf{그런데 $0.25$ 는 $0$ 밖 음수가 하나도 없고 앱이 $+0.0044$ 로
올랐다.} 사전 등록(``$0.5$ 가 봉우리'')이 반증됐다.

\textbf{씨앗 $40$ 으로 확인하니 뒤집혔다.}

\begin{center}
\begin{tabular}{lrr}
\hline
만화 $0.25$ & 씨앗 $20$ & 씨앗 $40$ \\
\hline
KR 만화 & $-0.0021$ & $\mathbf{-0.0035}$ $[-0.0060, -0.0009]$ \\
비게임 앱 & $+0.0044$ & $+0.0028$ $[-0.0003, +0.0059]$ \\
\hline
\end{tabular}
\end{center}

KR 이 $0$ 밖 음수가 되고 앱이 $0$ 을 문다. \textbf{기각 --- $0.5$ 유지.}
\end{abstract}

\begin{figure}[t]
\centering
\includegraphics[width=\textwidth]{figs/f1.pdf}
\caption{\textbf{가} --- 세 팔(씨앗 $20$). \textbf{나} --- $0.25$ 의
$20 \to 40$. \textbf{다} --- 문턱 곡선.}
\end{figure}

\section{씨앗 스물이 세 번째로 속였다}

\begin{center}
\begin{tabular}{lll}
\hline
후보 & 씨앗 $20$ & 씨앗 $40$ \\
\hline
세계애니 \texttt{grp}(노트 $504 \to 507$) & 판 $+0.0011$ $0$ 밖 & $+0.0000$ $0$ 을 뭄 → \textbf{되돌림} \\
만화 절반(노트 509) & 판 $+0.0020$ 경계 & $+0.0021$ 더 단단 → \textbf{채택} \\
만화 $0.25$(이 노트) & $0$ 밖 음수 없음 & KR $0$ 밖 음수 → \textbf{기각} \\
\hline
\end{tabular}
\end{center}

\textbf{세 번 중 두 번이 뒤집혔다.} 노트 507 이 만든 조항(판 차 $+0.002$
아래면 씨앗 $40$ 까지)이 그때마다 값을 했다. 이 노트에서는 판 차가
$-0.0000$ 이라 \emph{판으로는 아무 신호가 없었고}, 뒤집은 것은 집 밖
자였다.

\begin{quote}
\textbf{조항 넓히기} --- \textbf{씨앗 $40$ 문턱은 판 차뿐 아니라 어느
자에서든 결정이 $0$ 근처에 걸릴 때 적용한다.} 이 후보는 판이 완전히
비겨서 조항의 문구($+0.002$ 아래)에 걸렸는데, 실제로 판정을 뒤집은 것은
KR 이었다.
\end{quote}

\section{봉우리의 모양}

\begin{center}
\begin{tabular}{lrr}
\hline
만화 학습 행 비율 & 판 & KR 만화 \\
\hline
$0.25$ & $+0.0001$ & $\mathbf{-0.0035}$ \\
$\mathbf{0.5}$ & --- & --- \\
$0.75$ & $\mathbf{-0.0026}$ & $-0.0002$ \\
\hline
\end{tabular}
\end{center}

\textbf{양쪽이 다른 이유로 나쁘다} --- 너무 많이 남기면(0.75) 집안이
무너지고, 너무 많이 버리면(0.25) 집 밖이 무너진다. \textbf{$0.5$ 는 두
손해가 만나는 자리다.}

읽으면 --- 만화 행은 \emph{집안에는} 과하고 \emph{집 밖에는} 여전히
쓸모가 있다. 노트 508 이 ``만화는 받기만 하고 안 준다''고 했는데,
\textbf{집 밖에는 준다.}

\section*{이 사이클 --- 손잡이가 다 잠겼다}

\begin{center}
\begin{tabular}{lll}
\hline
손잡이 & 값 & 양옆을 쟀나 \\
\hline
\texttt{DOMDROP} & 애니 트렌드 & 노트 $497$(다른 도메인 · 다른 축) \\
\texttt{TRAINCAP} & 만화 $0.5$ & \textbf{이 노트}($0.25 \cdot 0.75$) \\
\texttt{TRAINW} & clip $(0.2, 1.0)$ & 노트 $510$(clip $2.0$ 기각) \\
\texttt{TRAINR} & $\tau = 8$ & 노트 $517$($4 \cdot 12 \cdot 16$) \\
\hline
\end{tabular}
\end{center}

\textbf{네 손잡이가 전부 양옆을 재고 잠겼다.} 판 $\mathbf{0.4932}$.

\section*{남은 것}

\begin{itemize}
\item \textbf{조합을 다시 열기} --- 넷을 하나씩 재느라 \emph{같이} 움직이는
경우를 안 봤다(예: $\tau$ 를 낮추면서 만화 문턱을 올리기).
\item \textbf{노트 446 거부권 다시 읽기}(노트 $516 \cdot 517$ 이 두 번
남겼다).
\item \textbf{자료 쪽} --- 노트 $505 \cdot 506$ 이 가리킨 곳. 이 사이클이
번 $+0.0120$ 은 ``행 두 배''의 절반이고, 나머지 절반은 \textbf{새 행이
있어야} 한다.
\end{itemize}

\end{document}
